\documentclass[10pt]{article}

\usepackage[a4paper,margin=0.82in]{geometry}
\usepackage{amsmath,amssymb,amsthm,mathtools,bm,mathrsfs}
\usepackage{booktabs,tabularx,microtype,enumitem,float}
\usepackage{xcolor}
\usepackage[colorlinks=true,linkcolor=blue!45!black,
  citecolor=blue!45!black,urlcolor=blue!45!black]{hyperref}
\setlength{\emergencystretch}{2em}
\setlist{nosep,leftmargin=*}

\newtheorem{theorem}{Theorem}
\newtheorem{corollary}[theorem]{Corollary}
\newtheorem{proposition}[theorem]{Proposition}
\newtheorem{lemma}[theorem]{Lemma}
\theoremstyle{definition}
\newtheorem{definition}[theorem]{Definition}
\newtheorem{remark}[theorem]{Remark}

\newcommand{\R}{\mathbb R}
\newcommand{\C}{\mathbb C}
\newcommand{\N}{\mathbb N}
\newcommand{\Z}{\mathbb Z}
\newcommand{\Q}{\mathbb Q}
\newcommand{\T}{\mathbb T}
\newcommand{\ee}{\mathrm e}
\newcommand{\ii}{\mathrm i}
\newcommand{\rank}{\operatorname{rank}}
\newcommand{\Rel}{\mathcal L}
\newcommand{\dist}{\operatorname{dist}}
\newcommand{\dens}{\operatorname{dens}}
\newcommand{\supp}{\operatorname{supp}}
\newcommand{\HS}{\mathrm{HS}}
\newcommand{\Hwater}{\mathcal H}
\newcommand{\proj}{\mathrm{proj}}

\title{\textbf{Persistent phase geometry of quantum recurrence}\\
\large\itshape Four structural laws for loss, spectral constraints, observability, and finite-horizon inference}
\author{GPT 5.6 Sol \and Opus 5 \and Fable 5 \and Michael Haufschild}
\date{24 August 2026}

\begin{document}
\maketitle

\begin{abstract}
Quantum recurrence is often described by one statement---finite isolated
systems eventually return---although loss, exact spectral resonances,
unequal state weights, restricted observables, and finite records answer
different questions.  We prove that these questions factor through one
persistent phase geometry.  First, for every finite-dimensional contraction,
peripheral spectral mass is the exact late-return ceiling: every lower
threshold recurs after every delay and every higher threshold eventually
fails.  Second, on the surviving phase carrier, balanced relation rank plus
independent projective phase rank is exactly $D-1$; the same lattice
determines the orbit hull, sharp generic first-hit exponents, and all excess
replica collisions.  Third, spectral weights pull a metric onto that hull.
A reverse-water functional controls both exact return-target information and
the least spectrum-uniform deadline within explicit constants, while any
finite probe family induces a coherence graph whose rank $D-c$ is the exact
observable phase capacity and the optimal universal recurrence exponent.
Fourth, completed rational rank is not identifiable from uniformly finite-
precision, finite-horizon data: arbitrarily close static l-bit Hamiltonians
can have maximal rank while their observed dynamics remain uniformly close
through the prescribed horizon.  The relation-hull and metric laws extend to
arbitrary countable pure-point packets with summable weights.

We then prove a physical saturation theorem absent from these abstract laws.
For every $n,q\geq2$, almost every static Hermitian onsite-plus-nearest-
neighbour Hamiltonian on every open simple-spectrum chart of a $q$-level
path has the maximum $q^n-1$ projective phase rank.  The same Hamiltonian has
sharp stroboscopic and first-gap-clock exponents $q^n-1$ and $q^n-2$, all-
order permutation-floor moments, and no nonzero exact projective return.
Its completed return density at chord tolerance $\varepsilon$ is exactly
$[2\arcsin(\varepsilon/2)/\pi]^{q^n-1}$ and therefore vanishes doubly
exponentially in chain length.  An exact commuting l-bit classification
identifies the contrasting boundary: finite interaction range can compress
phase rank, whereas spatial locality by itself cannot.
\end{abstract}

\noindent\textbf{Significance statement.}
Recurrence is governed neither by Hilbert-space dimension alone nor by a
single ``recurrence time.''  Loss first selects what survives; spectral
relations determine how many phases can move independently; state weights
set the accuracy demanded from each phase; probes decide which phase
differences are visible.  We give exact laws for all four reductions and
then show that ordinary static one-dimensional nearest-neighbour systems
generically saturate the largest possible phase carrier.  This turns a
qualitative inevitability into a structural and quantitative theory that
also explains when commuting locality, incomplete observation, or finite
precision makes the apparent recurrence problem smaller.

\section{One recurrence problem, four structural laws}

Quantum recurrence begins with a qualitative fact: a finite isolated system
returns arbitrarily close to its initial state \cite{BocchieriLoinger1957}.
Selected many-body recurrences have been observed directly, and coupled-
multiqubit experiments now target their size dependence
\cite{RauerEtAl2018,KarimiEtAl2026}.
That sentence suppresses four independent questions.  Does any part of the
initial condition survive indefinitely?  How many relative spectral phases
are genuinely independent?  How accurately must phases of unequal spectral
weight align?  Which of those alignments can a chosen experiment see?

Treating these questions as one ``recurrence time'' obscures their different
mathematics.  We instead factor recurrence through four objects:
\begin{equation}
 \boxed{
 \text{persistent carrier}
 \longrightarrow \text{relation hull}
 \longrightarrow \text{weighted metric}
 \longrightarrow \text{observer quotient}.}
 \label{eq:factorization}
\end{equation}
The carrier removes transient directions.  The balanced relation lattice
cuts the ambient phase torus down to the actual orbit closure.  State weights
pull a return metric onto that hull.  An observer graph finally quotients out
relative phases that no selected probe detects.  The first three laws
identify this factorization exactly, with the output of one reduction serving
as input to the next.  The fourth law identifies its finite-data boundary:
fixed-resolution metric statements are stable, but the completed arithmetic
relation lattice is not inferable from uniformly finite-precision,
finite-horizon records without an arithmetic promise.

\begin{table}[H]
\centering
\small
\begin{tabularx}{\textwidth}{@{}p{0.18\textwidth}p{0.24\textwidth}X@{}}
\toprule
Reduction & Exact object & What it controls \\
\midrule
Survival & unit-modulus spectral subspace & attainable late return amplitude \\
Constraints & balanced relation lattice & orbit hull, phase rank, exact resonances \\
Resolution & weighted projective metric & target volume and universal deadline \\
Observation & coherence-graph quotient & visible rank and probe-dependent return \\
\bottomrule
\end{tabularx}
\caption{Persistent phase geometry separates mechanisms that a single
dimension-only recurrence bound conflates.}
\label{tab:four-reductions}
\end{table}

The resulting four laws are followed by a physical saturation theorem.  A
generic $D\times D$ Hermitian matrix can have $D-1$ rationally independent
relative energies.  A one-dimensional nearest-neighbour Hamiltonian uses
only one-site and adjacent two-site terms.  It was therefore not known
whether this restricted coefficient space must carry hidden all-order energy
relations.  Nondegenerate gaps hold almost everywhere in broad local families
\cite{Huang2021}, but that excludes only pairwise gap coincidences.  Maximal
recurrence dimension excludes every balanced rational relation, at every
order.  The possible existence of hidden energy resonances in all two-body
Hamiltonians has been posed explicitly in autonomous state reachability
\cite{Szymanski2024}.

We prove that the full static open-path family has no such locality discount:
for every finite local dimension, almost every Hamiltonian on every simple
chart saturates the complete projective phase capacity.  The same draw then
obeys sharp first-hit laws, all-order collision extremality, and an exact
thermodynamic return-density law.  A complementary commuting l-bit theorem
shows where compression really enters: not through spatial locality alone,
but through exact algebraic restrictions on available interaction
coefficients \cite{SerbynPapicAbanin2013,HuseNandkishoreOganesyan2014,
VasseurParameswaranMoore2015}.

\subsection{Position relative to current results}

General recurrence bounds follow from simultaneous Diophantine approximation;
recent work obtains tight estimates and generic saturation for broad or
random finite systems \cite{KotowskiOszmaniec2026,GuptaShort2026}.  Weight-
sensitive concentration estimates and effective-support refinements address
related state-return problems \cite{RiddellPagliaroliAlhambra2023}.  Exact
Floquet recurrence has an arithmetic cyclotomic classification
\cite{AnandEtAl2026}.  Local free-fermion constructions can make their
single-particle energies rationally independent, while their many-body
spectra retain additive structure \cite{RiddellBertini2024}.  Spectral-form-
factor moments under independence conditions provide a complementary
collision viewpoint \cite{CamposVenutiOdavicHamma2026}.

Peripheral and attractor spaces are established tools for asymptotic open
dynamics \cite{WolfPerezGarcia2010,Albert2019,AmatoEtAl2023}; reverse water
filling is classical in rate--distortion theory
\cite{CoverThomas2006,GrayNeuhoff1998}; and Kronecker--Weyl, Dirichlet, and
metric extremality are established analytic ingredients
\cite{Weyl1916,KleinbockMargulis1998}.  Contributions here are the exact
joined factorization, weight- and observer-sensitive return laws, and the
local spectral saturation theorem that realizes their maximal regime in a
static bounded-degree one-dimensional many-body family.

\section{Law 1: loss selects the persistent carrier}
\label{sec:persistent-carrier}

Let $T$ be a linear contraction on a finite-dimensional complex Hilbert
space, $\|T\|\leq1$, and let $x$ be an initial vector.  Write
$\mathcal K_\partial$ for the span of the eigenspaces whose eigenvalues have
unit modulus, $\Pi_\partial$ for its orthogonal projection, and
\begin{equation}
 p_T(x)=\|\Pi_\partial x\|^2,
 \qquad A_n(x,T)=\langle x,T^nx\rangle .
 \label{eq:persistent-mass}
\end{equation}

\begin{theorem}[Law 1: exact returnability transition]
\label{thm:persistent-transition}
For every finite-dimensional contraction and every initial vector,
$0\leq p_T(x)\leq\|x\|^2$, and
\begin{align}
 a<p_T(x)&\Longrightarrow
   \forall M\ \exists n>M:\quad a\leq|A_n(x,T)|,          \label{eq:below-ceiling}\\
 p_T(x)<a&\Longrightarrow
   \exists B\ \forall n\geq B:\quad |A_n(x,T)|<a.         \label{eq:above-ceiling}
\end{align}
Thus $p_T(x)$ is the exact late-return-amplitude ceiling and $p_T(x)^2$ is
the exact fidelity ceiling.  If $p_T(x)>0$, the normalized weights
\begin{equation}
 w_\lambda=\frac{\|\Pi_\lambda x\|^2}{p_T(x)}
 \label{eq:peripheral-weights}
\end{equation}
form the finite phase packet governing every attainable threshold below the
ceiling.
\end{theorem}

This statement distinguishes decay from dephasing.  Components outside
$\mathcal K_\partial$ disappear; components inside it retain norm and differ
only by phase.  A unitary is the special case in which the full carrier
survives.  If the spectral radius is strictly below one, the persistent mass
is zero and every positive late-return threshold eventually becomes
impossible.  Nothing about later sections assumes weak loss: they apply to
the exact surviving packet selected by Theorem~\ref{thm:persistent-transition}.

\section{Law 2: spectral constraints determine the phase hull}
\label{sec:constraint-law}

Let $H$ be a finite-dimensional Hamiltonian with real energy packet
$E=(E_1,\ldots,E_D)$.  Global phase is physically irrelevant to recurrence,
so the natural characters are balanced.

\begin{definition}[Balanced relation lattice]
The projective spectral relation lattice is
\begin{equation}
 \Rel_E:=\left\{k\in\Z^D:
  \sum_{j=1}^D k_j=0,\quad \sum_{j=1}^D k_jE_j=0\right\}.
 \label{eq:lattice}
\end{equation}
Its rational rank is
$r(E):=\dim_\Q\operatorname{span}_\Q\Rel_E$.
For a reference level $E_1$, define the temporal rank
\begin{equation}
 \tau(E):=\dim_\Q\operatorname{span}_\Q
 \{E_j-E_1:2\leq j\leq D\}.
 \label{eq:temporal-rank}
\end{equation}

The spectrum is \emph{maximally projective} when $\Rel_E=\{0\}$,
equivalently when $\tau(E)=D-1$.
\end{definition}

The definition is independent of the chosen reference.  It includes all
finite-order no-resonance conditions at once: an equality between two sums
of the same number of energies is exactly an element of \eqref{eq:lattice}.

\begin{theorem}[Law 2: projective rank conservation]
For every nonempty finite spectrum,
\begin{equation}
 \boxed{\tau(E)+r(E)=D-1.}
 \label{eq:rank-conservation}
\end{equation}
More precisely, let
\begin{equation}
 \Phi_E:\{c\in\Q^D:\textstyle\sum_jc_j=0\}\longrightarrow\R,
 \qquad \Phi_E(c)=\sum_jc_jE_j.
 \label{eq:detuning-map}
\end{equation}

Then $\ker\Phi_E$ is the rational span of $\Rel_E$, and
$\operatorname{im}\Phi_E$ is the rational span of the anchored gaps.
\end{theorem}

\begin{proof}
The domain of $\Phi_E$ has rational dimension $D-1$.  Clearing
denominators shows that a rational balanced vector lies in the kernel exactly
when it lies in $\operatorname{span}_\Q\Rel_E$.  Every image element can be
written as
\[
 \sum_jc_jE_j=\sum_{j=2}^D c_j(E_j-E_1)
\]
because $\sum_jc_j=0$.  Conversely each anchored gap is the image of the
balanced vector $e_j-e_1$.  Rank--nullity gives
\eqref{eq:rank-conservation}.
\end{proof}

Equation \eqref{eq:rank-conservation} is the universal principle of the
paper.  An exact relation does not merely correlate two energies: it removes
one rational direction from the completed projective phase motion.  Maximal
projective rank means that the relative phase orbit is dense in the full
$(D-1)$-torus.  When $D\geq3$, it also rules out every nonzero \emph{exact}
simultaneous projective return; recurrence remains approximate rather than
periodic.

The lattice determines more than dimension.  Define its anchored form
\begin{equation}
 \Rel_E^{\rm anch}=
 \left\{\ell\in\Z^{D-1}:
 \left(-\sum_{j=1}^{D-1}\ell_j,\ell_1,\ldots,\ell_{D-1}\right)
 \in\Rel_E\right\}
 \label{eq:anchored-lattice}
\end{equation}
and the corresponding annihilator
\begin{equation}
 K_E=\left\{z\in\T^{D-1}:
 \prod_{j=1}^{D-1}z_j^{\ell_j}=1
 \text{ for every }\ell\in\Rel_E^{\rm anch}\right\}.
 \label{eq:relation-hull}
\end{equation}

\begin{theorem}[Exact relation-hull law]
The closure of the anchored phase line
\begin{equation}
 t\longmapsto
 \left(\ee^{-\ii t(E_2-E_1)},\ldots,
       \ee^{-\ii t(E_D-E_1)}\right)
 \label{eq:anchored-phase-line}
\end{equation}
is exactly $K_E$.  Consequently two finite spectra have the same completed
projective phase ensemble exactly when they have the same balanced relation
lattice.
\end{theorem}

Thus relation rank is the numerical shadow of a stronger object: the entire
compact group on which completed recurrence occurs.  Later metrics and
observers restrict this same $K_E$; they do not replace it.

\section{Sharp first-hit laws}

For $x\in\R$, write
$\lVert x\rVert_\Z:=\min_{p\in\Z}|x-p|$.  Given
$\beta=(\beta_1,\ldots,\beta_d)\in\R^d$ and $r>0$, define
\begin{equation}
 T_\beta(r):=\min\left\{q\in\N:q\geq1,
  \max_{1\leq j\leq d}\lVert q\beta_j\rVert_\Z<r\right\}.
 \label{eq:first-hit}
\end{equation}
Dirichlet's theorem gives $T_\beta(r)\leq\lceil1/r\rceil^d$.
The matching lower exponent holds outside the exceptional set of vectors
that are too well approximable.

\begin{definition}
A vector $\beta\in\R^d$ is very well approximable if some $\delta>0$
admits infinitely many $q\geq1$ for which
\[
 \max_j\lVert q\beta_j\rVert_\Z
 \leq q^{-1/d-\delta}.
\]
\end{definition}

\begin{theorem}[Metric first-hit exponent]
If $d\geq1$ and $\beta$ is not very well approximable, then
\begin{equation}
 \lim_{r\downarrow0}
 \frac{\log T_\beta(r)}{\log(1/r)}=d.
 \label{eq:first-hit-exponent}
\end{equation}
Consequently $T_\beta(r)=r^{-d+o(1)}$.
\end{theorem}

The metric theorem of Kleinbock and Margulis implies that almost every point
on a real-analytic nondegenerate manifold is not very well approximable
\cite{KleinbockMargulis1998}.  In the present setting, nondegeneracy is not
assumed.  Section~\ref{sec:local-proof} proves it from exact spectral escape
of the local Hamiltonian family.

For a simple spectrum, order energies as
$\lambda_1<\cdots<\lambda_D$ and introduce two canonical vectors:
\begin{align}
 \beta^{\rm disc}_j&=\frac{\lambda_{j+1}-\lambda_1}{2\pi},
   &&1\leq j\leq D-1,                                      \label{eq:disc-beta}\\
 \beta^{\rm wind}_j&=\frac{\lambda_{j+2}-\lambda_1}
                           {\lambda_2-\lambda_1},
   &&1\leq j\leq D-2.                                      \label{eq:wind-beta}
\end{align}
The first counts simultaneous anchored-gap returns at integer observation
times.  The second spends the first positive gap as a clock and counts the
winding needed to align every remaining gap.  It is a clock-normalized
winding statement, not a claim about the unrestricted continuous-time
first-return distribution.

\section{Replica collisions and return moments}

Let $w_j>0$ with $\sum_jw_j=1$.  These are spectral weights of a state with
nonzero overlap on every energy eigenvector.  Define survival amplitude and
fidelity
\begin{equation}
 A(t)=\sum_{j=1}^D w_j\ee^{-\ii E_jt},\qquad F(t)=|A(t)|^2.
 \label{eq:survival}
\end{equation}
For a multiplicity vector $p\in\N^D$, put
$|p|=\sum_jp_j$ and
$\binom{m}{p}=m!/\prod_jp_j!$.  The unavoidable order-$m$ permutation floor
is
\begin{equation}
 \Pi_m(w):=
 \sum_{\substack{p\in\N^D\\|p|=m}}
 \binom{m}{p}^{\!2}\prod_{j=1}^D w_j^{2p_j}.
 \label{eq:floor}
\end{equation}

\begin{theorem}[All-order collision equivalence]
For a finite spectrum and strictly positive weights, the symmetric long-time
moments satisfy
\begin{equation}
 \lim_{T\to\infty}\frac1{2T}\int_{-T}^{T}F(t)^m\,dt
 =\sum_{\substack{|p|=|q|=m\\p\cdot E=q\cdot E}}
 \binom{m}{p}\binom{m}{q}
 \prod_jw_j^{p_j+q_j}.                                    \label{eq:moment-collision}
\end{equation}
The right side is at least $\Pi_m(w)$.  Moreover, the following are
equivalent:
\begin{enumerate}[label=(\roman*)]
\item $\Rel_E=\{0\}$;
\item the collision sum on the right of \eqref{eq:moment-collision} equals
  $\Pi_m(w)$ for every $m\geq0$;
\item every long-time moment in \eqref{eq:moment-collision} equals
  $\Pi_m(w)$.
\end{enumerate}
The equivalence holds for every choice of strictly positive weights.
\end{theorem}

\begin{proof}
The multinomial expansion of $F(t)^m$ is a finite sum indexed by pairs
$(p,q)$ with $|p|=|q|=m$.  Its phase is
$\exp[-\ii t(p-q)\cdot E]$ and its coefficient is strictly positive.
Symmetric time averaging sends a nonzero detuning to zero and a zero
detuning to one, giving \eqref{eq:moment-collision}.  Terms with $p=q$ give
\eqref{eq:floor}.  If $\Rel_E=\{0\}$, every balanced equality
$(p-q)\cdot E=0$ forces $p=q$.  Conversely, any nonzero
$k\in\Rel_E$ has positive and negative parts $k^+,k^-$ with equal mass;
the unequal pair $(k^+,k^-)$ contributes a strictly positive excess at that
replica order.  This proves all equivalences.
\end{proof}

The theorem gives an operational reading of the full relation lattice.
Second-order gap nondegeneracy detects only a shallow part of it.  All
replica orders together are complete: every nontrivial balanced relation
appears at the order equal to its positive mass, and positivity prevents
cancellation from hiding it.

\section{Law 3: weights and observers set the visible return geometry}
\label{sec:weighted-observer}

The relation hull says which phase configurations are possible.  It does not
yet say how large a return target is.  Spectral weights supply that metric.
Let $w=(w_1,\ldots,w_r)$ have positive entries summing to one.  For a phase
packet $z\in\T^r$, define
\begin{equation}
 C_w(z)=\left|\sum_{j=1}^r w_jz_j\right|.
 \label{eq:weighted-coherence}
\end{equation}
The exact weighted variance identity
\begin{equation}
 1-C_w(z)^2=\frac12\sum_{i,j}w_iw_j|z_i-z_j|^2
 \label{eq:weighted-variance}
\end{equation}
shows that fidelity return is a projective phase-synchronization problem.

For a pivot $k$ and amplitude defect $0<\eta\leq1$, allocate distortions to
the non-pivot coordinates and set
\begin{align}
 \Hwater_k(w,\eta)
 &=\min\left\{\frac12\sum_{j\ne k}\log\frac{w_j}{D_j}:
 0<D_j\leq w_j,\ \sum_{j\ne k}D_j\leq\eta\right\},
 \label{eq:water-pivot}\\
 \Hwater_{\proj}(w,\eta)&=\min_k\Hwater_k(w,\eta).
 \label{eq:water-projective}
\end{align}
If $\eta<1-w_k$, its optimizer is
\begin{equation}
 D_j^*=\min\{w_j,\lambda_k^{-1}\},\qquad
 \sum_{j\ne k}\min\{w_j,\lambda_k^{-1}\}=\eta,
 \label{eq:reverse-water}
\end{equation}
and
\begin{equation}
 \Hwater_k(w,\eta)=\frac12\sum_{j\ne k}\log_+(\lambda_kw_j).
 \label{eq:water-value}
\end{equation}
This is numerically the classical parallel-Gaussian reverse-water program
\cite{CoverThomas2006,GrayNeuhoff1998}; here it describes deterministic phase
resolution, not a communication channel.

Let
\begin{align}
 \mu_w(\eta)&=m_{\T^r}\{z:C_w(z)\geq1-\eta\},             \label{eq:target-volume}\\
 I_w(\eta)&=-\log\mu_w(\eta),                              \label{eq:target-information}
\end{align}
and let $T_w(\eta)$ be the least $N$ such that every real phase packet
$\omega$ has some $1\leq n\leq N$ with
\begin{equation}
 \left|\sum_jw_j\ee^{\ii n\omega_j}\right|\geq1-\eta.
 \label{eq:universal-deadline}
\end{equation}

\begin{theorem}[Law 3a: weight--target--deadline law]
\label{thm:water-law}
For $r\geq2$ and $0<\eta\leq1$, the target has positive Haar volume and
\begin{align}
 \Hwater_{\proj}-\frac\pi4(r-1)-\log r
 &\leq I_w
 \leq\Hwater_{\proj}-(r-1)\log\frac{\sqrt2}{\pi},          \label{eq:target-sandwich}\\
 \Hwater_{\proj}-\frac\pi4(r-1)-\log(2r)
 &<\log T_w
 \leq\Hwater_{\proj}+(r-1)\log(2\pi\sqrt2).               \label{eq:deadline-sandwich}
\end{align}
For every phase packet and every delay $M\geq0$, a later return occurs at
some
\begin{equation}
 M<n\leq(M+1)T_w(\eta).                                    \label{eq:post-delay-return}
\end{equation}
Moreover, through the first $N$ ticks the Haar fraction of phase packets that
have returned is at most
\begin{equation}
 Nr\exp\left\{\frac\pi4(r-1)-\Hwater_{\proj}(w,\eta)\right\}.
 \label{eq:finite-horizon-haar}
\end{equation}
\end{theorem}

Hard support is therefore not the complete finite-resolution variable.
Small weights can fall below the water line and cease to contribute to the
leading synchronization cost.  The same functional controls the intrinsic
volume of the true target and the least deadline uniform over all compatible
spectra, but only within the displayed finite constants; no asymptotic
equality is being hidden.

Observation supplies a second quotient.  Let $E_1,\ldots,E_d$ be distinct
energy sectors and $B^{(a)}$ a finite family of states or observables with
weights $\lambda_a\geq0$.  Write $B^{(a)}_{ij}$ for energy-block matrix
elements and set
\begin{align}
 W_{ij}&=\sum_a\lambda_a\|B^{(a)}_{ij}\|_{\HS}^2,          \label{eq:observer-weight}\\
 \mathcal E_B(z)&=\sum_{i,j}W_{ij}|z_i-z_j|^2.             \label{eq:observer-energy}
\end{align}
The coherence graph $G_B$ has an edge $ij$ when $W_{ij}>0$.

\begin{theorem}[Law 3b: exact observer-coherence law]
\label{thm:observer-law}
For $U_t=\sum_i\ee^{-\ii tE_i}P_i$,
\begin{equation}
 \sum_a\lambda_a\|U_tB^{(a)}U_t^\dagger-B^{(a)}\|_{\HS}^2
 =\mathcal E_B(z(t)).                                      \label{eq:observer-identity}
\end{equation}
This vanishes exactly when phases agree on every edge, equivalently on every
connected component.  If $c(G_B)$ is the number of components, the exact
observable projective phase capacity is
\begin{equation}
 \boxed{d-c(G_B).}                                         \label{eq:observer-rank}
\end{equation}
For every integer $Q\geq1$, some
$1\leq n\leq(Q+1)^{d-c(G_B)}$ simultaneously places all visible phase
differences within the corresponding graph-error tolerance.  There are hard
spectra for which no smaller universal exponent is possible.
\end{theorem}

The graph theorem explains why global recurrence and observed recurrence can
have different effective dimensions without contradiction.  One common
phase per component is invisible.  Adding probes can connect components and
increase visible rank; changing the Hamiltonian is unnecessary.

\section{Law 4: finite horizons reveal geometry, not completed arithmetic}
\label{sec:finite-horizon}

Exact rational rank is an infinite-precision invariant.  Finite-time unitary
dynamics is norm-continuous, while the relation lattice is arithmetically
discontinuous.  This produces a sharp division between what finite records
can certify and what completed recurrence knows.

\begin{theorem}[Law 4: finite-horizon arithmetic ambiguity]
\label{thm:finite-horizon-ambiguity}
For every finite binary l-bit Hamiltonian $H_J$, every $\varepsilon>0$, and
every finite observation horizon $T>0$, there is a Hamiltonian
$H_{\widetilde J}$ diagonal in the same product basis such that
\begin{align}
 \sum_S|\widetilde J_S-J_S|&<\varepsilon,                  \label{eq:lbit-close}\\
 \tau(E_{\widetilde J})&=2^N-1,                            \label{eq:lbit-maximal-close}
\end{align}
while all off-diagonal configuration transitions remain zero.  For every
normalized observable $O$ and $|t|\leq T$,
\begin{equation}
 \|\ee^{-\ii tH_{\widetilde J}}O\ee^{\ii tH_{\widetilde J}}
   -\ee^{-\ii tH_J}O\ee^{\ii tH_J}\|
 \leq2T\varepsilon\|O\|.                                 \label{eq:finite-window-close}
\end{equation}
Thus no uniformly finite-precision, finite-horizon record can determine the
completed exact relation rank without an additional arithmetic promise.
\end{theorem}

This no-go does not make finite-time recurrence uninformative.  Equations
\eqref{eq:deadline-sandwich} and \eqref{eq:finite-horizon-haar} quantify
metric statements that are stable at fixed resolution.  The distinction is
between completed arithmetic equality and finite-resolution target geometry.

The carrier law is not confined to a finite spectrum.  Let $(E_j)_{j\geq0}$
be a countable pure-point packet with weights $p_j\geq0$, $\sum_jp_j=1$.
Anchor $\omega_j=E_j-E_0$, let $\Rel_E^{\rm fin}$ be the lattice of finitely
supported balanced integer relations, and let $K_E$ be its annihilator in the
countable torus.  The weighted phase embedding
\begin{equation}
 \Psi_p(z)=(\sqrt{p_j}z_j)_{j\geq0}\in\ell^2
 \label{eq:weighted-embedding}
\end{equation}
pulls the Hilbert metric onto the exact projective hull.

\begin{theorem}[Countable relation--metric extension]
\label{thm:countable-extension}
The closure of the anchored countable phase line is exactly $K_E$, and
anchoring identifies its finite-character relation lattice with
$\Rel_E^{\rm fin}$.  For every $r>0$, a positive survival return with squared
defect at most $2r^2$ occurs no later than the covering number
\begin{equation}
 N_r\bigl(\Psi_p(K_E)\bigr).                               \label{eq:countable-cover}
\end{equation}
In particular, every normalized summable pure-point packet recurs at every
positive weighted phase resolution, after every finite checkpoint.
\end{theorem}

Finite rank is therefore not the universal primitive.  The complete
relation hull is; finite rank is its finite-dimensional shadow, and covering
entropy replaces dimension for countably infinite pure-point carriers.

\section{Static nearest-neighbour paths attain the maximum}
\label{sec:local-main}

Let
\(
 \mathscr H_{n,q}=(\C^q)^{\otimes n}
\)
and let $\mathscr V_{n,q}$ be the real vector space of all Hermitian operators
that are sums of terms supported on one site or one edge of the open path:
\begin{equation}
 \mathscr V_{n,q}=
 \left\{\sum_{j=1}^{n}A_j^{(j)}+
              \sum_{j=1}^{n-1}B_j^{(j,j+1)}:
 A_j=A_j^\dagger,\ B_j=B_j^\dagger\right\}.
 \label{eq:qudit-path-space}
\end{equation}
Superscripts denote tensor embedding.  Fix once and for all a real basis
$P_1,\ldots,P_N$ of this space and write
\begin{equation}
 H_{n,q}(\theta)=\sum_{a=1}^{N}\theta_aP_a,
 \qquad \theta\in\R^N.
 \label{eq:qudit-path-family}
\end{equation}
The identity direction is included; it shifts every energy equally and hence
drops out of all balanced relations.  The interaction graph is the open path,
so its maximum degree is two, and no longer-range term is present.

\begin{theorem}[Same-Hamiltonian local saturation]
\label{thm:local-capstone}
Let $n,q\geq2$, $D=q^n$, and let $U\subset\R^N$ be any nonempty,
connected, open set on which $H_{n,q}(\theta)$ has simple spectrum.  For
Lebesgue-almost every $\theta\in U$, one and the same Hamiltonian satisfies:
\begin{enumerate}[label=(\alph*)]
\item $\Rel_{E(\theta)}=\{0\}$ and $\tau(E(\theta))=D-1$;
\item for the ordered spectrum, the stroboscopic vector
  \eqref{eq:disc-beta} obeys
  \[
   \lim_{r\downarrow0}
   \frac{\log T_{\beta^{\rm disc}}(r)}{\log(1/r)}=D-1;
  \]
\item the clock-normalized vector \eqref{eq:wind-beta} obeys
  \[
   \lim_{r\downarrow0}
   \frac{\log T_{\beta^{\rm wind}}(r)}{\log(1/r)}=D-2;
  \]
\item for every strictly positive spectral weight and every $m\geq0$, the
  replica collision mass and the long-time survival-fidelity moment equal
  the permutation floor \eqref{eq:floor}.
\end{enumerate}

In addition, there is no nonzero exact simultaneous projective return.
\end{theorem}

The key point is quantifier order.  A single full-measure set of coefficients
works simultaneously for maximal rank, both exponent laws, every positive
weight, and every replica order.  The Hamiltonian is not changed between
clauses.  A nonempty open simple-spectrum chart exists for every $n,q\geq2$.

\begin{corollary}[No locality discount]
For almost every Hamiltonian in each chart of Theorem~\ref{thm:local-capstone},
the stroboscopic return exponent is $q^n-1$.  Thus spatial dimension one,
bounded interaction degree, and only nearest-neighbour terms do not reduce
the exponent below the full projective Hilbert-space dimension.
\end{corollary}

\section{The exact commuting-local boundary}
\label{sec:lbit-boundary}

The local saturation theorem does not say that every local family is
maximal.  It says that geometric locality alone cannot force a deficit.  A
commuting l-bit family makes the missing algebraic restriction explicit.
For $N$ binary modes and $x\in\{0,1\}^N$, let
\begin{equation}
 E_J(x)=\sum_{S\subseteq\{1,\ldots,N\}}J_S
       \prod_{i\in S}(-1)^{x_i}.                           \label{eq:lbit-energy}
\end{equation}
All terms commute and the configuration basis is fixed.

\begin{theorem}[Exact l-bit phase-rank classification]
\label{thm:lbit-classification}
For every coefficient family,
\begin{equation}
 \boxed{
 \tau(E_J)=\dim_\Q\operatorname{span}_\Q
 \{J_S:S\neq\varnothing\}.}                              \label{eq:lbit-rank}
\end{equation}
Consequently
\begin{equation}
 \tau(E_J)+r(E_J)=2^N-1.                                  \label{eq:lbit-conservation}
\end{equation}
If only one-body coefficients are present, $\tau(E_J)\leq N$.  If nonzero
supports have diameter at most $R$, then
\begin{equation}
 \tau(E_J)\leq N,2^{R+1}.                                \label{eq:lbit-range}
\end{equation}
Both support ceilings are attained by rationally independent allowed
coefficients, while rational independence of all nonconstant $J_S$ gives
$\tau(E_J)=2^N-1$.
\end{theorem}

Walsh orthogonality recovers every $J_S$ from the complete energy packet, so
coefficient span and gap span coincide.  The contrast with
Theorem~\ref{thm:local-capstone} is exact.  Fixed-range commuting support
permits only linearly many coefficient directions at fixed $R$; the full
noncommuting nearest-neighbour path, despite also using linearly many local
parameters, generically maps those parameters into $q^n-1$ rationally
independent eigenvalue directions.  Parameter-count dimension and rational
spectral rank are different invariants.

One concrete state sees the entire l-bit law.  The uniform product state
$|+\rangle^{\otimes N}$ has positive weight on every configuration, so its
visible gap space equals the complete gap space.  Its all-order long-time
survival moments attain the permutation floors exactly when
$\tau(E_J)=2^N-1$.  The rank is therefore not merely hidden in an
unobservable enumeration of eigenvalues.

\section{Exact completed density and the thermodynamic law}

Fix a reference energy $E_1$ and define the projective chord defect
\begin{equation}
 \Delta_E(t)=\max_{2\leq j\leq D}
 \left|\ee^{-\ii t(E_j-E_1)}-1\right|.
 \label{eq:chord-defect}
\end{equation}
\par\smallskip\noindent
For a measurable time set $S$, use centered completed density
\begin{equation}
 \dens(S)=\lim_{T\to\infty}
 \frac{\bigl|S\cap[-T,T]\bigr|}{2T},
 \label{eq:density}
\end{equation}
\par\smallskip\noindent
whenever the limit exists.

\begin{theorem}[Exact projective-return density]
If $E$ is maximally projective and $0<\varepsilon<2$, then
\begin{equation}
 \boxed{
 \dens\{t:\Delta_E(t)<\varepsilon\}
 =\left(\frac{2\arcsin(\varepsilon/2)}{\pi}\right)^{D-1}.}
 \label{eq:exact-density}
\end{equation}
The same value is obtained with a non-strict inequality.
\end{theorem}

\begin{proof}
Maximal projective rank makes the anchored phase flow equidistributed on
$\T^{D-1}$.  In one circle, chord distance below $\varepsilon$ from the
identity is an arc of normalized Haar measure
$2\arcsin(\varepsilon/2)/\pi$.  The target is the Cartesian product of
$D-1$ such arcs.  Its boundary has Haar measure zero, so equidistribution
gives \eqref{eq:exact-density} for strict and closed targets.
\end{proof}

\subsection{Thermodynamic limit}

For an $n$-site $q$-level path, set
\begin{align}
 b_\varepsilon=\frac{2\arcsin(\varepsilon/2)}{\pi},\qquad
 I_\varepsilon&=-\log b_\varepsilon>0,                       \label{eq:rate}\\
 \rho_{n,q}(\varepsilon)&=b_\varepsilon^{q^n-1}.             \label{eq:rho-definition}
\end{align}

\begin{theorem}[Thermodynamic projective-recurrence extinction]
For every $q\geq2$ and $0<\varepsilon<2$, the exact finite-size identity and
its three limits are
\begin{align}
 -\log\rho_{n,q}(\varepsilon)&=(q^n-1)I_\varepsilon,          \label{eq:finite-rate}\\
 \rho_{n,q}(\varepsilon)&\longrightarrow0,                  \label{eq:rho-zero}\\
 \frac{-\log\rho_{n,q}(\varepsilon)}{q^n}
   &\longrightarrow I_\varepsilon,                          \label{eq:hilbert-rate}\\
 \frac{\log[-\log\rho_{n,q}(\varepsilon)]}{n}
   &\longrightarrow\log q.                                  \label{eq:double-log}
\end{align}
\end{theorem}

\paragraph{Order of limits.}
Equation \eqref{eq:double-log} is the precise double-exponential statement.
The order of limits matters: completed-time density is taken at each fixed
finite $n$, and only then is $n\to\infty$ taken.  This theorem does not
construct or assert recurrence of an infinite-volume global orbit.

\section{Comparison with the literature}

Existing results illuminate individual layers of
\eqref{eq:factorization}.  Bocchieri and Loinger prove qualitative finite
quantum recurrence \cite{BocchieriLoinger1957}.  General and random-model
work gives recurrence bounds, escape-time dependence, concentration, and
generic saturation without the exact local-family relation theorem
\cite{RiddellPagliaroliAlhambra2023,KotowskiOszmaniec2026,GuptaShort2026}.
Peripheral spectra and attractor algebras classify asymptotic channel
structure \cite{WolfPerezGarcia2010,Albert2019,AmatoEtAl2023}; our contraction
law isolates the state-dependent persistent mass as a sharp scalar threshold
and attaches a finite weighted deadline below it.  Reverse water filling is
classical \cite{CoverThomas2006,GrayNeuhoff1998}; the recurrence result is that
the same finite optimization controls both the exact projective target volume
and the least spectrum-uniform return deadline, with explicit two-sided
remainders.  The coherence-graph law supplies the observer quotient that a
state-only entropy does not encode.

For locality, Huang proves nondegenerate gaps for almost every geometrically
local Hamiltonian, but pairwise gap nondegeneracy is weaker than excluding
every balanced relation \cite{Huang2021}.  Riddell and Bertini construct local
free-fermion models with rationally independent single-particle spectra, while
their many-body spectra retain additive structure \cite{RiddellBertini2024}.
Exact spectral-form-factor moments under independence hypotheses are prior
art \cite{CamposVenutiOdavicHamma2026}; our local theorem supplies the missing
all-order independence inside the full interacting path family.  Exact
Floquet recurrence has a complementary cyclotomic arithmetic classification
\cite{AnandEtAl2026}.

To our knowledge, none of the cited results proves the joint statement of
Theorem~\ref{thm:local-capstone}: almost-everywhere maximal projective rank,
the $D-1$ and $D-2$ sharp first-hit exponents, and all-order moment floors for
one Hamiltonian drawn from the full static nearest-neighbour qudit path
family.  In particular, the theorem gives a negative answer, for this
canonical family and modulo global energy shift, to the proposed obstruction
that every two-body Hamiltonian might carry a hidden energy resonance
\cite{Szymanski2024}.

\section{Interpretation and scope}

The factorization changes what ``recurrence complexity'' means.  Loss can
remove amplitude without changing the relation geometry of what survives.
Relations can shrink the completed orbit without changing hard spectral
support.  Small weights can remove directions from finite-resolution cost
without changing exact rank.  Restricted probes can identify only a quotient
of the same phase hull.  These are distinct operations and should not be
compressed into one effective dimension.

The local result then isolates a common misconception.  Locality restricts
the support and number of Hamiltonian coefficients, but projective recurrence
rank is the rational rank of the resulting eigenvalue gaps.  A low-dimensional
analytic coefficient family can be arithmetically nondegenerate in a much
larger spectral space.  Conversely, fixed-range commuting l-bit structure
really does bound rank because Walsh duality identifies spectral directions
one-for-one with allowed interaction coefficients.

The scope is exact and deliberately bounded.
\begin{itemize}
\item Loss law concerns finite-dimensional contractions.  It does not claim
  that every quantum channel is Hilbert--Schmidt contractive.
\item Hamiltonian laws concern finite, closed, time-independent systems;
  countable extension concerns normalized summable pure-point phase packets.
\item Generic means Lebesgue-almost every coefficient in the stated local
  family or chart, not every local Hamiltonian.  Commuting and integrable
  subfamilies can have many exact relations.
\item Exponent $D-1$ is for integer-time anchored-gap hits.  Exponent $D-2$
  is for first-gap-clock-normalized winding.  Neither is presented as an
  unrestricted continuous-time first-return distribution.
\item Exact density is a completed-time statement.  It gives no finite-time
  convergence rate without additional quantitative Diophantine information.
\item Thermodynamic limit is ordered: time completion precedes growth of
  system size.  No infinite-volume global recurrence is claimed.
\item Finite-horizon ambiguity concerns exact arithmetic inference.  It does
  not preclude stable estimation of metric, finite-resolution quantities.
\end{itemize}

Within those boundaries, the conclusion is end-to-end: persistence selects
the phase carrier; relations determine its exact hull; weights and probes set
its operational geometry; and a static one-dimensional bounded-degree family
generically realizes the largest finite projective carrier allowed by its
Hilbert space.

\appendix

\section{Proof of the persistent-carrier law}
\label{app:persistent}

Unit-modulus eigenspaces of a finite contraction are mutually orthogonal and
reduce both $T$ and $T^\dagger$.  Decompose
\begin{equation}
 x=\sum_{\lambda\in\Lambda_\partial}\Pi_\lambda x+x_{\rm tr}.
 \label{eq:peripheral-decomposition}
\end{equation}
The restriction of $T$ to the orthogonal complement of
$\mathcal K_\partial$ has spectral radius strictly below one, hence its
powers converge to zero in operator norm.  Therefore
\begin{equation}
 A_n(x,T)=\sum_{\lambda\in\Lambda_\partial}
 \lambda^n\|\Pi_\lambda x\|^2+o(1).                       \label{eq:peripheral-sum}
\end{equation}
The nonnegative coefficients in the phase sum total $p_T(x)$.  Simultaneous
phase recurrence makes the sum arbitrarily close to this total after every
delay, while the triangle inequality bounds its modulus by the same total.
The vanishing transient gives \eqref{eq:below-ceiling} and
\eqref{eq:above-ceiling}.  Normalization gives
\eqref{eq:peripheral-weights}.

\section{Proof of the relation-hull and weighted laws}
\label{app:hull-water}

For the anchored phase line, a character $\ell\in\Z^{D-1}$ is constant one
exactly when
\begin{equation}
 \sum_{j=1}^{D-1}\ell_j(E_{j+1}-E_1)=0,
 \label{eq:anchored-character-zero}
\end{equation}
which is equivalent to the balanced $D$-coordinate relation in
\eqref{eq:anchored-lattice}.  The character-annihilator form of
Kronecker--Weyl therefore identifies the orbit closure with
\eqref{eq:relation-hull}.  Since a compact subgroup is determined by its
annihilating characters, equality of relation lattices and equality of
completed phase ensembles are equivalent.

For the water optimizer, put $x_j=w_j/D_j$.  Feasibility becomes
$x_j\geq1$ and $\sum_{j\ne k}w_j/x_j\leq\eta$, while the objective is
$\frac12\sum_{j\ne k}\log x_j$.  The candidate
$x_j^*=\max\{1,\lambda_kw_j\}$ satisfies the constraint.  Summing the
one-coordinate supporting inequality
\begin{equation}
 \log\max\{1,\lambda a\}
 +\lambda\left(\min\{a,\lambda^{-1}\}-\frac ax\right)
 \leq\log x,qquad x\geq1,                                \label{eq:water-support}
\end{equation}
proves global optimality and \eqref{eq:water-value}.

For the target lower bound, use a minimizing pivot and require the projective
coordinates to lie in arcs with
\begin{equation}
 \|y_j\|_\T\leq\frac{\sqrt{2/x_j}}{2\pi}.
 \label{eq:water-box}
\end{equation}
The weighted variance identity places this product box inside the target.
Its exact normalized Haar volume is
\begin{equation}
 \left(\frac{\sqrt2}{\pi}\right)^{r-1}
 \exp\{-\Hwater_{\proj}\}.                                \label{eq:water-box-volume}
\end{equation}
For the upper bound, resolve a returned packet by a pivot of maximal weighted
alignment, apply an exponential envelope to each phase coordinate, and use
the optimized water level.  This gives
\begin{equation}
 \mu_w(\eta)\leq
 r\exp\left\{\frac\pi4(r-1)-\Hwater_{\proj}\right\},       \label{eq:water-upper-volume}
\end{equation}
which together with \eqref{eq:water-box-volume} proves
\eqref{eq:target-sandwich}.

For the constructive deadline, partition each non-pivot circle into
\begin{equation}
 Q_j=\left\lceil\pi\sqrt{2x_j}\right\rceil
 \label{eq:water-grid}
\end{equation}
arcs.  Two among the first $\prod_{j\ne k}Q_j+1$ orbit points lie in the
same product box; their time difference returns.  Bounding the rounded
product yields the upper half of \eqref{eq:deadline-sandwich}.  Conversely,
every power map preserves Haar measure.  The union of the first $N$ return
targets has measure at most $N$ times
\eqref{eq:water-upper-volume}; choosing $N$ below its reciprocal yields a
phase packet with no earlier return and proves the lower half.  Applying the
deadline to $(M+1)\omega$ proves \eqref{eq:post-delay-return}.  The same union
bound proves \eqref{eq:finite-horizon-haar}.

\section{Proof of the observer, finite-horizon, and countable laws}
\label{app:observer-countable}

In the energy basis,
\begin{equation}
 (U_tB^{(a)}U_t^\dagger-B^{(a)})_{ij}
 =(z_i\overline z_j-1)B^{(a)}_{ij}.
 \label{eq:block-phase}
\end{equation}
Summing squared Hilbert--Schmidt block norms proves
\eqref{eq:observer-identity}.  Every summand is nonnegative, so the sum
vanishes exactly when phases agree across every positive edge.  The kernel
of the graph coboundary consists of functions constant on connected
components; rank--nullity gives $d-c(G_B)$.  Synchronizing one root-relative
phase per coboundary direction by Dirichlet approximation gives the stated
$(Q+1)^{d-c(G_B)}$ ceiling.  Choosing rationally independent hard frequencies
on a graph basis gives the matching minimax exponent.

For Theorem~\ref{thm:finite-horizon-ambiguity}, the nonconstant coefficient
packet lies in $\R^{2^N-1}$.  Rationally dependent packets form the countable
union of rational hyperplanes, so every open $\ell^1$ ball contains a
rationally independent packet $\widetilde J$.  Walsh inversion then gives
\eqref{eq:lbit-maximal-close}.  Both Hamiltonians remain diagonal in the same
configuration basis.  The Duhamel identity and
$\|H_{\widetilde J}-H_J\|\leq\sum_S|\widetilde J_S-J_S|$ give
\eqref{eq:finite-window-close}.

For the countable extension, finitely supported characters separate points
of the countable torus.  Balancing a character by moving its total charge to
the anchor preserves detuning, so the anchored and balanced relation
lattices coincide.  The character-annihilator theorem identifies the orbit
closure with $K_E$.  Moreover,
\begin{equation}
 \|\Psi_p(z)-\Psi_p(w)\|_{\ell^2}^2
 =\sum_jp_j|z_j-w_j|^2,                                   \label{eq:pullback-metric}
\end{equation}
so an $r$-cover of $\Psi_p(K_E)$ pulls back to a finite cover of the exact
phase hull.  Two orbit points in one cover element have squared separation at
most $4r^2$, hence weighted coherence defect at most $2r^2$, which gives
\eqref{eq:countable-cover}.  Translating the finite cover along the orbit
gives a hit after every checkpoint.  Compactness follows from the summable
tail and finite-head approximation.

\section{Proof of the l-bit classification}
\label{app:lbit}

The Walsh characters form an orthogonal basis on the Boolean cube.  Every
energy gap in \eqref{eq:lbit-energy} is a rational linear combination of the
nonconstant $J_S$, while Walsh inversion expresses each $J_S$ as a rational
linear combination of energy values and hence of an anchored energy and
gaps.  The coefficient span and projective gap span are equal, proving
\eqref{eq:lbit-rank}; Law~2 gives \eqref{eq:lbit-conservation}.

If only one-body supports occur, at most $N$ coefficients generate the span.
For diameter at most $R$, assign each nonempty support to its leftmost site.
It then lies in a window of at most $R+1$ sites, giving at most $2^{R+1}$
choices per left endpoint and hence \eqref{eq:lbit-range}.  Rationally
independent allowed coefficients attain the corresponding support count.
Finally, the uniform product state has equal positive weight on every Boolean
configuration, so its visible gap span is the complete one; the collision
equivalence proved in the main text gives the moment statement.

\section{Proof of the first-hit exponent}
\label{app:first-hit}

We record the short argument behind Law~2.  Dirichlet approximation gives
\begin{equation}
 T_\beta(r)\leq\lceil1/r\rceil^d,
 \label{eq:dirichlet-upper}
\end{equation}
and therefore
\[
 \limsup_{r\downarrow0}
 \frac{\log T_\beta(r)}{\log(1/r)}\leq d.
\]
Suppose $\beta$ is not very well approximable.  For each $c>1/d$, only
finitely many $q$ satisfy
$\max_j\lVert q\beta_j\rVert_\Z\leq q^{-c}$.  Hence, below a sufficiently
small radius, a first hit $q=T_\beta(r)$ must satisfy $q^{-c}<r$, or
$q>r^{-1/c}$.  Thus
\[
 \liminf_{r\downarrow0}
 \frac{\log T_\beta(r)}{\log(1/r)}\geq\frac1c.
\]
Letting $c\downarrow1/d$ gives the lower bound $d$, matching
\eqref{eq:dirichlet-upper}.

\section{Local spectral escape and genericity}
\label{sec:local-proof}

We give the algebraic mechanism behind Theorem~\ref{thm:local-capstone}.
For an unlabeled spectrum, escape must hold for every relabeling.

\begin{definition}[Affine spectral escape]
A real Hermitian family $\mathscr V$ has affine spectral escape if, for every
nonzero $c\in\R^D$ with $\sum_i c_i=0$ and every $a\in\R$, there are
$H\in\mathscr V$ and distinct real roots $E_1,\ldots,E_D$ of its
characteristic polynomial such that
\begin{equation}
 \sum_i c_iE_{\sigma(i)}\neq a
 \qquad\text{for every }\sigma\in S_D.                     \label{eq:aff-escape}
\end{equation}
\end{definition}

The load-bearing step is a one-leaf rigidity statement.  Let $E_\alpha$ be
a maximally projective old spectrum, let $r$ label a new $q$-level factor,
and let $A$ be a Hermitian endpoint transition matrix in the old energy
basis.  Write $G_A$ for the graph whose edge $\alpha\beta$ is present when
$A_{\alpha\beta}A_{\beta\alpha}\neq0$.  Its cross-layer lift has vertices
$(\alpha,r)$ and edges
\begin{equation}
 (\alpha,r)\sim(\beta,s)
 \quad\Longleftrightarrow\quad
 \alpha\sim_{G_A}\beta\ \text{ and }\ r\neq s.             \label{eq:lift}
\end{equation}

\begin{lemma}[Leaf-extension affine escape]
\label{lem:leaf-escape}
Suppose the graph in \eqref{eq:lift} is connected.  Choose real shifts
$s_r$ for which all $E_\alpha+s_r$ are distinct.  Then the real Hermitian
slice
\begin{equation}
 M(u,B)=\operatorname{diag}(E_\alpha+u_r)+A\otimes B,
 \qquad u\in\R^q,\quad B=B^{\mathsf T}\in M_q(\R),          \label{eq:leaf-slice}
\end{equation}
has affine spectral escape.
\end{lemma}

\begin{proof}
Near $(u,B)=(s,0)$, simplicity gives analytic eigenvalue branches indexed by
$(\alpha,r)$.  Assume a coefficient packet $c_{\alpha r}$ has constant
spectral action on an open neighborhood.  Fix distinct layers $r,s$, vary
their relative onsite shift through infinitely many collision-free values,
and restrict $B$ to the $rs$ exchange direction.  The second derivative in
the exchange strength gives a finite simple-pole identity.  Maximal
projectivity of $E$ makes the oriented old-energy gaps distinct.  Clearing
denominators and evaluating at each pole therefore gives
\begin{equation}
 c_{\alpha r}=c_{\beta s}
 \quad\text{whenever}\quad
 (\alpha,r)\sim(\beta,s).                                  \label{eq:edge-equality}
\end{equation}
Connectivity of the lifted graph makes every coefficient equal.  If $c$ is
balanced, that common value is zero.  Thus no nonzero balanced spectral
action can be constant, regardless of its proposed target $a$.

Apply this argument to each of the finitely many permuted coefficient
packets.  Each corresponding analytic action is nonconstant on the connected
chart, so its prescribed level set is nowhere dense.  Their finite union
cannot fill the chart.  A point outside that union gives
\eqref{eq:aff-escape}, with simple real eigenvalues by construction.
\end{proof}

The connectedness hypothesis is available uniformly in local dimension.

\begin{lemma}[Interaction-lift connectivity]
If $G$ is connected and nonbipartite and $q\geq2$, then the graph on
$V(G)\times\{1,\ldots,q\}$ defined by \eqref{eq:lift} is connected.
\end{lemma}

\begin{proof}
For $q=2$, walking along an edge flips the layer label.  Connectedness moves
between old vertices, while an odd closed walk, which exists because $G$ is
nonbipartite, changes parity and joins the two labels above one vertex.  For
$q\geq3$, a third label routes around any unwanted equality of endpoint
labels; ordinary connected paths in $G$ then lift between arbitrary pairs.
\end{proof}

\begin{proposition}[Path affine spectral escape]
\label{prop:path-escape}
For every $n,q\geq2$, the exact path-local space $\mathscr V_{n,q}$ in
\eqref{eq:qudit-path-space} has affine spectral escape.
\end{proposition}

\begin{proof}
For $n=2$, the single edge supports both sites, so
$\mathscr V_{2,q}$ is the full Hermitian matrix space.  A real diagonal
matrix with distinct entries can avoid the finitely many affine hyperplanes
in \eqref{eq:aff-escape}.

For a longer path, use its first $n-1$ sites as the old system and attach the
last site as one leaf.  The required old-system witness is obtained by a
finite leaf schedule.  If $q\geq3$, begin with a one-site simple maximally
projective spectrum and a complete one-site Hermitian probe family.  If
$q=2$, begin with a two-site simple maximally projective seed whose endpoint
probe family has connected nonbipartite transition support.  At each later
attachment, the leaf-extension construction preserves maximal projectivity
and supplies a connected nonbipartite one-site probe family at every retained
site.  Thus the endpoint of the $(n-1)$-site path has a Hermitian linear
combination $A$ with connected nonbipartite support.

Choose integer-coded layer shifts so that the product energies
$E_\alpha+s_r$ are all distinct.  The interaction-lift lemma and
Lemma~\ref{lem:leaf-escape} show that the new-onsite-plus-last-edge slice has
affine spectral escape.  Conjugating from the old energy basis back to the
physical tensor basis turns \eqref{eq:leaf-slice} into
\begin{equation}
 K\otimes I+I\otimes\operatorname{diag}(u)+A_{\rm end}\otimes B.
\end{equation}
Here $K$ is local on the first $n-1$ sites, $A_{\rm end}$ is supported only
on their endpoint, and the last term is supported only on the new edge.
Hence the entire slice lies in $\mathscr V_{n,q}$, proving the proposition.
\end{proof}

\begin{lemma}[Genericity on every simple chart]
\label{lem:chart-genericity}
For the family \eqref{eq:qudit-path-family}:
\begin{enumerate}[label=(\roman*)]
\item Hamiltonians with a nontrivial balanced integer spectral relation form
  a Lebesgue-null set;
\item on every connected open simple-spectrum chart, both maps
  $\theta\mapsto\beta^{\rm disc}(\theta)$ and
  $\theta\mapsto\beta^{\rm wind}(\theta)$ are real analytic and
  nondegenerate in the metric-Diophantine sense.
\end{enumerate}
\end{lemma}

\begin{proof}
For a fixed nonzero balanced $k\in\Z^D$, form the orbit product
\begin{equation}
 \mathcal O_k(E)=\prod_{\sigma\in S_D}
                 \left(\sum_i k_iE_{\sigma(i)}\right).      \label{eq:orbit-product}
\end{equation}
It is symmetric in the energies, hence a polynomial in the characteristic
coefficients, which are polynomial in $\theta$.  Proposition
\ref{prop:path-escape} supplies a point where $\mathcal O_k\neq0$; thus its
pullback is a nonzero real polynomial.  Its zero set is null.  The countable
union over balanced $k\neq0$ proves (i).

On a simple chart, ordered eigenvalues are real analytic.  A real affine
relation among the coordinates of either gap map, constant on a nonempty
open subset, extends throughout the connected chart.  For the discrete gap
map it is directly a balanced affine spectral action.  For the ratio map,
multiplication by the nonzero first gap produces the same kind of action.
Symmetrizing over labelings would contradict Proposition
\ref{prop:path-escape}.  Hence $1$ and the coordinate functions of either map
are linearly independent on every open subset, which is analytic
nondegeneracy.
\end{proof}

\begin{proof}[Proof of Theorem~\ref{thm:local-capstone}]
Lemma~\ref{lem:chart-genericity}(i) gives maximal projective rank almost
everywhere.  Metric extremality of nondegenerate analytic maps
\cite{KleinbockMargulis1998}, followed by Law~2, gives exponents $D-1$ and
$D-2$ for the two gap maps.  Intersecting these three conull events leaves a
single conull event.  On it, Law~3 gives every collision and moment conclusion
for all positive weights and all orders.  Finally $D=q^n\geq4$; maximal
projectivity makes the anchored frequencies rationally independent, so a
common exact phase return implies $t=0$.
\end{proof}

\section{Thermodynamic identities}

For completeness, $0<b_\varepsilon<1$ follows from
$0<\arcsin(\varepsilon/2)<\pi/2$.  Equation
\eqref{eq:finite-rate} is the logarithm of
$\rho_{n,q}=b_\varepsilon^{q^n-1}$.  Dividing by $q^n$ gives
$(1-q^{-n})I_\varepsilon\to I_\varepsilon$, proving
\eqref{eq:hilbert-rate}.  Since $q^n-1\to\infty$ and
$b_\varepsilon<1$, \eqref{eq:rho-zero} follows.  Finally,
\[
 \frac{\log[-\log\rho_{n,q}]}n
 =\frac{\log(q^n-1)+\log I_\varepsilon}{n}
 \longrightarrow\log q,
\]
which proves \eqref{eq:double-log}.


\section*{Author contributions}
All authors contributed to conception, mathematical development, analysis,
verification, and writing.

\begin{thebibliography}{99}

\bibitem{BocchieriLoinger1957}
P. Bocchieri and A. Loinger,
Quantum recurrence theorem,
\emph{Phys. Rev.} \textbf{107}, 337--338 (1957),
\href{https://doi.org/10.1103/PhysRev.107.337}{doi:10.1103/PhysRev.107.337}.

\bibitem{Weyl1916}
H. Weyl,
\"Uber die Gleichverteilung von Zahlen mod. Eins,
\emph{Math. Ann.} \textbf{77}, 313--352 (1916),
\href{https://doi.org/10.1007/BF01475864}{doi:10.1007/BF01475864}.

\bibitem{KleinbockMargulis1998}
D. Y. Kleinbock and G. A. Margulis,
Flows on homogeneous spaces and Diophantine approximation on manifolds,
\emph{Ann. Math.} \textbf{148}, 339--360 (1998),
\href{https://doi.org/10.2307/120997}{doi:10.2307/120997}.

\bibitem{RiddellPagliaroliAlhambra2023}
J. Riddell, N. J. Pagliaroli, and \'A. M. Alhambra,
Concentration of quantum equilibration and an estimate of the recurrence time,
\emph{SciPost Phys.} \textbf{15}, 165 (2023),
\href{https://doi.org/10.21468/SciPostPhys.15.4.165}{doi:10.21468/SciPostPhys.15.4.165}.

\bibitem{Huang2021}
Y. Huang,
Extensive entropy from unitary evolution,
arXiv:2104.02053 (2021),
\href{https://doi.org/10.48550/arXiv.2104.02053}{doi:10.48550/arXiv.2104.02053}.

\bibitem{RiddellBertini2024}
J. Riddell and B. Bertini,
Rationally independent free fermions with local hopping,
\emph{Phys. Rev. E} \textbf{110}, 064101 (2024),
\href{https://doi.org/10.1103/PhysRevE.110.064101}{doi:10.1103/PhysRevE.110.064101}.

\bibitem{KotowskiOszmaniec2026}
M. Kotowski and M. Oszmaniec,
Tight bounds on recurrence time in closed quantum systems,
\emph{Phys. Rev. Lett.} \textbf{137}, 060402 (2026),
\href{https://doi.org/10.1103/crwj-qfw5}{doi:10.1103/crwj-qfw5}.

\bibitem{GuptaShort2026}
C. Gupta and A. J. Short,
Recurrence time for finite quantum systems,
arXiv:2604.14995 (2026),
\href{https://doi.org/10.48550/arXiv.2604.14995}{doi:10.48550/arXiv.2604.14995}.

\bibitem{CamposVenutiOdavicHamma2026}
L. Campos Venuti, J. Odavi\'c, and A. Hamma,
Integrability and chaos via fractal analysis of spectral form factors:
Gaussian approximations and exact results,
\emph{Phys. Rev. E} \textbf{113}, 044128 (2026),
\href{https://doi.org/10.1103/5xw1-x4kh}{doi:10.1103/5xw1-x4kh}.

\bibitem{AnandEtAl2026}
A. Anand, D. Valluri, J. Davis, and S. Ghose,
Quantum recurrences and the arithmetic of Floquet dynamics,
\emph{Quantum} \textbf{10}, 2074 (2026),
\href{https://doi.org/10.22331/q-2026-04-20-2074}{doi:10.22331/q-2026-04-20-2074}.

\bibitem{Szymanski2024}
K. Szyma\'nski,
State reachability in isolated systems with 2-body Hamiltonians,
presentation (2024),
\href{https://quantumstat.es/talks/29aug2024/}{quantumstat.es/talks/29aug2024/}.

\bibitem{RauerEtAl2018}
B. Rauer, S. Erne, T. Schweigler, F. Cataldini, M. Tajik, and
J. Schmiedmayer,
Recurrences in an isolated quantum many-body system,
\emph{Science} \textbf{360}, 307--310 (2018),
\href{https://doi.org/10.1126/science.aan7938}{doi:10.1126/science.aan7938}.

\bibitem{KarimiEtAl2026}
B. Karimi, X. Wu, A. N. Cleland, and J. P. Pekola,
Blueprint for experiments exploring the quantum recurrence theorem on a
coupled multiqubit system,
\emph{Phys. Rev. Research} \textbf{8}, L012062 (2026),
\href{https://doi.org/10.1103/wv9h-x14y}{doi:10.1103/wv9h-x14y}.

\bibitem{WolfPerezGarcia2010}
M. M. Wolf and D. P\'erez-Garc\'ia,
The inverse eigenvalue problem for quantum channels,
arXiv:1005.4545 (2010),
\href{https://arxiv.org/abs/1005.4545}{arXiv:1005.4545}.

\bibitem{Albert2019}
V. V. Albert,
Asymptotics of quantum channels: conserved quantities, an adiabatic limit,
and matrix product states,
\emph{Quantum} \textbf{3}, 151 (2019),
\href{https://doi.org/10.22331/q-2019-06-06-151}{doi:10.22331/q-2019-06-06-151}.

\bibitem{AmatoEtAl2023}
D. Amato, P. Facchi, and A. Konderak,
Asymptotics of quantum channels,
\emph{J. Phys. A: Math. Theor.} \textbf{56}, 265304 (2023),
\href{https://doi.org/10.1088/1751-8121/acd828}{doi:10.1088/1751-8121/acd828}.

\bibitem{CoverThomas2006}
T. M. Cover and J. A. Thomas,
\emph{Elements of Information Theory}, 2nd ed. (Wiley, 2006).

\bibitem{GrayNeuhoff1998}
R. M. Gray and D. L. Neuhoff,
Quantization,
\emph{IEEE Trans. Inf. Theory} \textbf{44}, 2325--2383 (1998),
\href{https://doi.org/10.1109/18.720541}{doi:10.1109/18.720541}.

\bibitem{SerbynPapicAbanin2013}
M. Serbyn, Z. Papi\'c, and D. A. Abanin,
Local conservation laws and the structure of the many-body localized states,
\emph{Phys. Rev. Lett.} \textbf{111}, 127201 (2013),
\href{https://doi.org/10.1103/PhysRevLett.111.127201}{doi:10.1103/PhysRevLett.111.127201}.

\bibitem{HuseNandkishoreOganesyan2014}
D. A. Huse, R. Nandkishore, and V. Oganesyan,
Phenomenology of fully many-body-localized systems,
\emph{Phys. Rev. B} \textbf{90}, 174202 (2014),
\href{https://doi.org/10.1103/PhysRevB.90.174202}{doi:10.1103/PhysRevB.90.174202}.

\bibitem{VasseurParameswaranMoore2015}
R. Vasseur, S. A. Parameswaran, and J. E. Moore,
Quantum revivals and many-body localization,
\emph{Phys. Rev. B} \textbf{91}, 140202(R) (2015),
\href{https://doi.org/10.1103/PhysRevB.91.140202}{doi:10.1103/PhysRevB.91.140202}.

\end{thebibliography}


\end{document}
